\section{Sponsored scalar finality}
\label{sec:scalar-finality}

We now forget all physical witnesses.  The only data retained from Theorem~\ref{thm:uniform-profile} are the scalar demand $d_\rho(Q)$ and, at prime currents, the scalar mint $m_\rho(q)$.  The purpose of this section is to prove target terminality before any actual Haxell packing is chosen.

Enumerate only the late current subtype
\[
\{Q=p^e:Q\ge Q_{\mathrm{prof}}\}
\]
in increasing rank.  Let a grade interval $I=[A,B]$ have width
\[
w(I)=B-A.
\]
For a prime event $q$, Theorem~\ref{thm:uniform-profile} gives the action
\[
[A,B]\longmapsto
[A+d_\rho(q),\,B+d_\rho(q)+m_\rho(q)],
\tag{6.1}\label{eq:prime-action}
\]
which increases width by $m_\rho(q)$.  For a composite event $Q=p^e$, $e\ge2$, it gives
\[
[A,B]\longmapsto[A+d_\rho(Q),B],
\tag{6.2}\label{eq:composite-action}
\]
which decreases width by $d_\rho(Q)$.  An event is legal whenever its incoming width is at least its demand; in that case the new interval overlaps the old one.

Put
\[
\Phi(X)=\frac X{\log X}.
\]
Let
\[
A(X)=\max_{Q\le X}d_\rho(Q)
\]
over the late event domain, and let
\[
N(X)=
\sum_{\substack{p^e\le X\\ e\ge2}}d_\rho(p^e)
\tag{6.3}\label{eq:cumulative-debt}
\]
be the cumulative negative variation.  By \eqref{eq:d-composite},
\[
A(X)=O_\rho((\log X)^2)=o(\Phi(X)).
\tag{6.4}\label{eq:max-demand-subscale}
\]
There are $O(\sqrt X\log X)$ prime powers $p^e\le X$ with $e\ge2$: sum the crude bounds $X^{1/e}$ for $2\le e\le\log_2X$.  Hence
\[
\boxed{
N(X)=O_\rho(\sqrt X\,\log^3X)
=o(\Phi(X)).}
\tag{6.5}\label{eq:debt-subscale}
\]

\subsection{Earlier prime sponsors}

Let $\Lambda\ge2$ be the integer multiplier from Proposition~\ref{prop:prime-supply-main}.  There is a fixed integer
\[
\Lambda_s=4\Lambda
\tag{6.6}\label{eq:sponsor-lag}
\]
such that every sufficiently large event rank $Q$ has an earlier prime sponsor $q$ satisfying
\[
q<Q\le\Lambda_s q.
\tag{6.7}\label{eq:sponsor-bounded-lag}
\]
Indeed, put
\[
Y=Q\operatorname{div}(2\Lambda).
\]
For $Q$ sufficiently large, $Y$ is in the range of Proposition~\ref{prop:prime-supply-main}, so there is a prime
\[
Y<q\le\Lambda Y.
\]
Euclidean division gives
\[
Q=(2\Lambda)Y+r,
\qquad 0\le r<2\Lambda.
\]
Thus $q\le\Lambda Y<2\Lambda Y\le Q$, while for $Y\ge1$,
\[
Q<2\Lambda(Y+1)\le4\Lambda Y<4\Lambda q.
\]
Hence \eqref{eq:sponsor-bounded-lag}.  Enlarging the onset so that $Q>\Lambda_sQ_{\mathrm{prof}}$ also forces $q>Q_{\mathrm{prof}}$, so the sponsor belongs to the same late event subtype.

By \eqref{eq:mint-prime} and \eqref{eq:sponsor-bounded-lag},
\[
m_\rho(q)
\gg_\rho \frac q{\log q}
\asymp_{\rho,\Lambda_s}\frac Q{\log Q}
=\Theta(\Phi(Q)).
\tag{6.8}\label{eq:sponsor-mint}
\]
Thus every sufficiently late event has a bounded-multiplicative-lag earlier prime sponsor whose positive width increment is macroscopic on the scale $\Phi(Q)$.

\subsection{Startup reserve and legal continuation}

The scalar theorem will be used above an arbitrarily prescribed lower cutoff $B_\dagger$.  Define the startup request
\[
R(B)=A(\Lambda_sB)+N(\Lambda_sB).
\tag{6.9}\label{eq:startup-request}
\]
By \eqref{eq:max-demand-subscale}, \eqref{eq:debt-subscale}, and fixed-scale stability of $\Phi$,
\[
R(B)=o(B/\log B).
\tag{6.10}
\]
Hence Proposition~\ref{prop:wide-start} applies to this particular request.  Choose $B$ sufficiently large so that
\[
B\ge B_\dagger,
\qquad
B\ge Q_{\mathrm{prof}},
\]
and so that all asymptotic statements above and the WideStart conclusion hold.  Fix one resulting WideStart with grade interval of width at least $R(B)$ and full centre $\gamma_B$ from \eqref{eq:full-centre}.

If an event $Q$ has its chosen sponsor $q$ at or before the start $B$, then \eqref{eq:sponsor-bounded-lag} implies
\[
Q\le\Lambda_sB.
\]
Pessimistically discard every positive mint in this finite transition.  The initial width $R(B)$ covers both the largest individual demand and all negative variation through $\Lambda_sB$.

After every sponsor lies inside the post-start chain, retain only the positive increment of one sponsor and subtract all negative variation through the current event.  Immediately before an event of rank $Q$ we have the lower bound
\[
w(Q)\ge m_\rho(q)-N(Q).
\tag{6.11}\label{eq:sponsored-width}
\]
By \eqref{eq:sponsor-mint} and \eqref{eq:debt-subscale},
\[
w(Q)\gg\Phi(Q),
\]
whereas the current demand is $o(\Phi(Q))$ by \eqref{eq:max-demand-subscale}.  Thus every sufficiently late event is legal.  Together with the startup reserve, induction gives legality of the entire post-start chain.

Successive integer intervals overlap, and prime sponsor mints are unbounded because their ranks are cofinal and $\Phi(X)\to\infty$.  Therefore the union of the scalar grade intervals contains a final ray in $\mathbf N$.

\subsection{The centre dies at a finite endpoint}

The WideStart begins in the affine fibre centred at the fixed class
\[
\gamma_B=\alpha_B+\beta_B\in A/E_B.
\]
Every productive current response after the start is internal to the endpoint filtration and transports this centre by the canonical quotient map.  By cofinality of $(E_X)$, there exists a finite cutoff $X_{\mathrm{abs}}$ after which the image of $\gamma_B$ in $A/E_X$ is zero.  Deleting the finite initial part of an interval chain does not change the final-ray conclusion.

For a grade $k$ define
\[
\operatorname{Term}_k=
\{X\ge B:
\gamma_B\mapsto0\text{ in }A/E_X,
\quad k\in I_X\},
\tag{6.12}\label{eq:term-k}
\]
where $I_X$ is the scalar grade interval after processing all late currents through $X$.

\begin{theorem}[Scalar terminality]\label{thm:scalar-terminality}
For every prescribed lower cutoff $B_\dagger$, one can choose a WideStart base $B\ge B_\dagger$ such that
\[
\exists k_0\ \forall k\ge k_0:
\qquad \operatorname{Term}_k\ne\varnothing.
\]
No Haxell packing, conflict graph, physical union, or reciprocal-resource tail is used in this target-only theorem.
\end{theorem}
